\scanpage{104}

Galois motivique de $\overline{\mathbb F}_p$. \deleted{\unsure{Comme les} $M^\circ(q^{\unsure{m}})\quad M^\circ(q)$ \unsure{induisent un isom.}
\[
M^\circ(q')\otimes_ZQ\quad M^\circ(q)\otimes_ZQ,
\]}
on constate que $M^\circ(q^\infty)$ \deleted{que} s'identifie à \deleted{\unsure{un sous-groupe}} de $M^\circ\deleted{(\unsure{q})}\otimes_ZQ$, qui est lui-même isomorphe au quotient $P/U$, $U$ étant le groupe des racines de l'1, et $P$ étant le sous-groupe de $\bar Q^*$ formé des $\lambda$ tels que $\lambda^n\in M^\circ\deleted{(\unsure{q})}$ pour $n$ assez grand.

\margin{On aura $M(q^\infty)\simeq Z\times M^\circ(q^\infty)$, $M^\circ(q^\infty)=\varinjlim M^\circ(q^n)\simeq M^\circ\otimes_ZQ$.}

\deleted{L'hom
\diagram{\begin{tikzcd}[column sep=large]
M^\circ(q^{\deleted{\unsure{m}}}=p^n)\arrow[r]&P/U
\end{tikzcd}}
s'obtient par
\[
\lambda\longrightarrow\lambda^{1/n}\bmod U
\]
(où la détermination de $\unsure{U|\lambda^{1/n}}$ qu'on prend est sans importance). On a bien
\[
(\lambda^{1/n})^{\unsure{na}}=\lambda^{\unsure{a}}\in M^\circ(p^n)^{\unsure{a}}
\]
\inserted{$M^{\unsure{a}}(p)$}
car $M(p^n)^{\unsure{a}}\subset M(p)$. Trouve-t-on $P^\circ/U$ tout entier? \deleted{\unsure{On peut}} \inserted{Si} $\lambda\in\bar Q^*$ tels que $\lambda^n\in M^\circ(p)$, \deleted{\unsure{donc si}} provient-il du $\mu\in M^\circ(p^m)$, i.e. \deleted{\unsure{est il de}} a-t-on $\lambda^m\in M^\circ(p^m)$ pour $m$ convenable??

On peut \deleted{\unsure{supposer de $\lambda$ s'exprime}} se borner aux $m=nh$ multiples de $n$, et la question devient si pour tt $\lambda\in M^\circ(p)$, $n\deleted{\unsure{h}}\in N^*$ étant donnés, il existe un $h\in N^*$ tels que $\lambda^h\in M^\circ(p^{hn})$.}

\mspage{103}
